\documentclass[12pt,a4paper]{article}

% ─── packages ─────────────────────────────────────────────────────────────────
\usepackage[a4paper, margin=2.5cm]{geometry}
\usepackage{amsmath, amssymb, amsfonts}
\usepackage{graphicx}
\usepackage{booktabs}
\usepackage{array}
\usepackage{xcolor}
\usepackage{listings}
\usepackage{caption}
\usepackage{subcaption}
\usepackage{hyperref}
\usepackage{titlesec}
\usepackage{enumitem}
\usepackage{fancyhdr}
\usepackage{float}
\usepackage{mathtools}

% ─── colors ───────────────────────────────────────────────────────────────────
\definecolor{codeblue}{RGB}{30,90,160}
\definecolor{codegray}{RGB}{100,100,100}
\definecolor{codegreen}{RGB}{0,130,80}
\definecolor{backcolor}{RGB}{245,246,248}
\definecolor{accentblue}{RGB}{15,52,96}

% ─── listings style ───────────────────────────────────────────────────────────
\lstdefinestyle{cppstyle}{
    backgroundcolor=\color{backcolor},
    commentstyle=\color{codegreen}\itshape,
    keywordstyle=\color{codeblue}\bfseries,
    stringstyle=\color{orange},
    basicstyle=\ttfamily\footnotesize,
    breakatwhitespace=false,
    breaklines=true,
    captionpos=b,
    keepspaces=true,
    numbers=left,
    numberstyle=\tiny\color{codegray},
    showspaces=false,
    showstringspaces=false,
    showtabs=false,
    tabsize=4,
    language=C++,
    frame=single,
    rulecolor=\color{codegray!50},
    xleftmargin=1.5em,
    xrightmargin=0.5em,
    aboveskip=0.8em,
    belowskip=0.8em
}
\lstset{style=cppstyle}

% ─── header/footer ────────────────────────────────────────────────────────────
\pagestyle{fancy}
\fancyhf{}
\lhead{\small Kalman Filter — Milestone 2}
\rhead{\small Full-Body 3D Gait Estimation}
\cfoot{\thepage}
\renewcommand{\headrulewidth}{0.4pt}

% ─── section formatting ───────────────────────────────────────────────────────
\titleformat{\section}{\large\bfseries\color{accentblue}}{\thesection.}{0.6em}{}[\titlerule]
\titleformat{\subsection}{\normalsize\bfseries\color{accentblue!80}}{\thesubsection.}{0.5em}{}

% ─── hyperref ─────────────────────────────────────────────────────────────────
\hypersetup{
    colorlinks=true,
    linkcolor=accentblue,
    urlcolor=codeblue,
    pdftitle={Kalman Filter Milestone 2},
}

% ─── custom commands ──────────────────────────────────────────────────────────
\newcommand{\bx}{\hat{\mathbf{x}}}
\newcommand{\bF}{\mathbf{F}}
\newcommand{\bQ}{\mathbf{Q}}
\newcommand{\bH}{\mathbf{H}}
\newcommand{\bP}{\mathbf{P}}
\newcommand{\bK}{\mathbf{K}}
\newcommand{\bR}{\mathbf{R}}
\newcommand{\bI}{\mathbf{I}}
\newcommand{\bz}{\mathbf{z}}
\newcommand{\by}{\mathbf{y}}
\newcommand{\bJ}{\mathbf{J}}

% ══════════════════════════════════════════════════════════════════════════════
\begin{document}

% ─── title page ───────────────────────────────────────────────────────────────
\begin{titlepage}
    \centering
    \vspace*{2cm}
    {\Huge\bfseries\color{accentblue} Kalman Filter\\[0.4em]
    \large Milestone 2\par}
    \vspace{1.5cm}
    {\Large Linear and Extended Kalman Filter\\
    Implementation for 3D Full-Body Human Gait Estimation\par}
    \vspace{2cm}
    \rule{0.6\textwidth}{0.4pt}\\[0.5cm]
    {\large C/C++ Implementation $\cdot$ Full-Body State Estimation\\
    Numerical Design $\cdot$ Plot Analysis $\cdot$ 3D Simulation\par}
    \vspace{1.5cm}
    \rule{0.6\textwidth}{0.4pt}\\[0.5cm]
    {\large\today}
    \vfill
    \begin{abstract}
    \noindent
    This report presents the full implementation of a Linear Kalman Filter (LKF) and
    Extended Kalman Filter (EKF) for estimating the 3D kinematic state of a 23-joint
    human body during walking. Each joint is modelled with a 12-dimensional constant-jerk
    state vector, yielding a 276-dimensional full-body state. Both filters are implemented
    in C/C++ with careful attention to numerical stability, computational efficiency, and
    memory management. The LKF uses a linear Cartesian measurement model while the EKF
    employs a nonlinear spherical coordinate model requiring a manually implemented
    arctan2 approximation. Filtered outputs are analysed through time-series plots, RMSE
    comparisons, and 3D full-body walking animations.
    \end{abstract}
\end{titlepage}

% ─── table of contents ────────────────────────────────────────────────────────
\tableofcontents
\newpage

% ══════════════════════════════════════════════════════════════════════════════
\section{Introduction}

Tracking human motion from noisy sensor data is a fundamental problem in biomechanics,
rehabilitation engineering, and sports science. Motion capture systems provide position
measurements that are inherently corrupted by sensor noise, making direct use of raw data
unreliable for applications requiring velocity, acceleration, or higher-order derivatives.

The Kalman Filter provides an optimal recursive estimator for linear Gaussian systems,
combining a physics-based prediction step with a measurement correction step to yield
minimum-variance estimates. Where the measurement model is nonlinear—such as when
observations are expressed in spherical coordinates—the Extended Kalman Filter linearises
the nonlinear function at each timestep using its Jacobian.

This milestone implements both filters in C/C++ for a 23-joint full-body gait dataset,
producing smooth estimates of position, velocity, acceleration, and jerk for all joints.

% ══════════════════════════════════════════════════════════════════════════════
\section{Dataset Description}

\subsection{Source}
The dataset is a 3D Full Body Human Gait Walking Dataset containing motion data for
23 body joints. Two CSV files are provided:
\begin{itemize}[noitemsep]
    \item \texttt{true\_values.csv} — ground truth Cartesian positions for all joints
    \item \texttt{noisy\_values.csv} — position measurements corrupted with additive Gaussian noise
\end{itemize}

\subsection{Structure}
Each CSV contains 3{,}040 rows (frames) and 69 columns representing the $x$, $y$, $z$
coordinates of 23 joints. The joint names are: pelvis, L5, L3, T12, T8, neck, head,
shoulderRight, upperArmRight, forearmRight, handRight, shoulderLeft, upperArmLeft,
forearmLeft, handLeft, upperLegRight, lowerLegRight, footRight, toeRight, upperLegLeft,
lowerLegLeft, footLeft, and toeLeft.

\subsection{Sampling Rate}
The dataset corresponds to a motion capture system operating at 100~Hz, giving a
sampling interval of:
\[
    \Delta t = \frac{1}{100} = 0.01 \ \text{s}
\]

% ══════════════════════════════════════════════════════════════════════════════
\section{State-Space Model}

\subsection{Single-Joint State Vector}
For one joint, the 12-dimensional state vector captures position, velocity,
acceleration, and jerk along all three Cartesian axes:
\[
    \mathbf{x}_k = \begin{bmatrix}
        p_x & v_x & a_x & j_x &
        p_y & v_y & a_y & j_y &
        p_z & v_z & a_z & j_z
    \end{bmatrix}^\top \in \mathbb{R}^{12}
\]

\subsection{Full-Body State Vector}
For 23 independent joints, states are stacked:
\[
    \mathbf{x}_{\text{body}} = \begin{bmatrix}
        \mathbf{x}_{\text{joint}_1}^\top &
        \mathbf{x}_{\text{joint}_2}^\top &
        \cdots &
        \mathbf{x}_{\text{joint}_{23}}^\top
    \end{bmatrix}^\top \in \mathbb{R}^{276}
\]

\subsection{Constant-Jerk Dynamics}
The motion model assumes jerk is approximately constant over one sampling interval:
\[
    \dot{p} = v, \quad \dot{v} = a, \quad \dot{a} = j, \quad \dot{j} = 0
\]
Discretising via Taylor expansion yields:
\begin{align}
    p_{k+1} &= p_k + v_k\Delta t + \tfrac{1}{2}a_k\Delta t^2 + \tfrac{1}{6}j_k\Delta t^3 \\
    v_{k+1} &= v_k + a_k\Delta t + \tfrac{1}{2}j_k\Delta t^2 \\
    a_{k+1} &= a_k + j_k\Delta t \\
    j_{k+1} &= j_k
\end{align}

% ══════════════════════════════════════════════════════════════════════════════
\section{System Matrices}

\subsection{State Transition Matrix \texorpdfstring{$\mathbf{F}$}{F}}
For one axis, the $4\times4$ block is:
\[
    F_{\text{axis}} = \begin{bmatrix}
        1 & \Delta t & \frac{\Delta t^2}{2} & \frac{\Delta t^3}{6} \\
        0 & 1        & \Delta t             & \frac{\Delta t^2}{2} \\
        0 & 0        & 1                    & \Delta t             \\
        0 & 0        & 0                    & 1
    \end{bmatrix}
\]
The full $12\times12$ matrix is block-diagonal over three independent axes:
\[
    \bF = \begin{bmatrix} F_{\text{axis}} & \mathbf{0} & \mathbf{0} \\
                          \mathbf{0} & F_{\text{axis}} & \mathbf{0} \\
                          \mathbf{0} & \mathbf{0} & F_{\text{axis}} \end{bmatrix}
\]

\subsection{Process Noise Covariance \texorpdfstring{$\mathbf{Q}$}{Q}}
Noise enters through jerk: $j_{k+1} = j_k + \eta_k$, $\eta_k \sim \mathcal{N}(0,\sigma_j^2)$.
The propagation vector is $\mathbf{g} = [\Delta t^3/6,\ \Delta t^2/2,\ \Delta t,\ 1]^\top$.
Then $Q_{\text{axis}} = \sigma_j^2\, \mathbf{g}\mathbf{g}^\top$, expanded as:
\[
    Q_{\text{axis}} = \sigma_j^2
    \begin{bmatrix}
        \frac{\Delta t^6}{36} & \frac{\Delta t^5}{12} & \frac{\Delta t^4}{6}  & \frac{\Delta t^3}{6}  \\
        \frac{\Delta t^5}{12} & \frac{\Delta t^4}{4}  & \frac{\Delta t^3}{2}  & \frac{\Delta t^2}{2}  \\
        \frac{\Delta t^4}{6}  & \frac{\Delta t^3}{2}  & \Delta t^2            & \Delta t              \\
        \frac{\Delta t^3}{6}  & \frac{\Delta t^2}{2}  & \Delta t              & 1
    \end{bmatrix}
\]
Full $\bQ$ is block-diagonal with $\sigma_j^2 = 1.0$ and $\Delta t = 0.01$.

\subsection{Measurement Matrix \texorpdfstring{$\mathbf{H}$}{H} (LKF)}
The sensor measures only Cartesian position:
\[
    \bH = \begin{bmatrix}
        1 & 0 & 0 & 0 & 0 & 0 & 0 & 0 & 0 & 0 & 0 & 0 \\
        0 & 0 & 0 & 0 & 1 & 0 & 0 & 0 & 0 & 0 & 0 & 0 \\
        0 & 0 & 0 & 0 & 0 & 0 & 0 & 0 & 1 & 0 & 0 & 0
    \end{bmatrix} \in \mathbb{R}^{3\times12}
\]

\subsection{EKF Nonlinear Measurement Model}
For the EKF, measurements are expressed in spherical coordinates:
\[
    h(\mathbf{x}) = \begin{bmatrix} r \\ \theta \\ \phi \end{bmatrix} =
    \begin{bmatrix}
        \sqrt{p_x^2 + p_y^2 + p_z^2} \\
        \text{atan2}(p_y,\, p_x) \\
        \text{atan2}\!\left(p_z,\, \sqrt{p_x^2+p_y^2}\right)
    \end{bmatrix}
\]

\subsection{EKF Jacobian \texorpdfstring{$\mathbf{J}$}{J}}
Let $\rho = \sqrt{p_x^2+p_y^2}$ and $r = \sqrt{p_x^2+p_y^2+p_z^2}$.
The $3\times12$ Jacobian (non-zero entries only in columns 1, 5, 9) is:
\[
    \bJ = \begin{bmatrix}
        \frac{p_x}{r}         & \cdots & \frac{p_y}{r}         & \cdots & \frac{p_z}{r}         & \cdots \\
        -\frac{p_y}{\rho^2}   & \cdots & \frac{p_x}{\rho^2}    & \cdots & 0                     & \cdots \\
        -\frac{p_x p_z}{\rho r^2} & \cdots & -\frac{p_y p_z}{\rho r^2} & \cdots & \frac{\rho}{r^2} & \cdots
    \end{bmatrix}
\]
where $\cdots$ represents zero columns for $v_x, a_x, j_x, \ldots$

% ══════════════════════════════════════════════════════════════════════════════
\section{Linear Kalman Filter Implementation}

\subsection{Filter Equations}
\textbf{Prediction:}
\begin{align}
    \hat{\mathbf{x}}_{k|k-1} &= \bF\hat{\mathbf{x}}_{k-1|k-1} \\
    \bP_{k|k-1}               &= \bF\bP_{k-1|k-1}\bF^\top + \bQ
\end{align}
\textbf{Update:}
\begin{align}
    \by_k     &= \bz_k - \bH\hat{\mathbf{x}}_{k|k-1} \\
    \bK_k     &= \bP_{k|k-1}\bH^\top\left(\bH\bP_{k|k-1}\bH^\top + \bR\right)^{-1} \\
    \hat{\mathbf{x}}_{k|k} &= \hat{\mathbf{x}}_{k|k-1} + \bK_k\by_k \\
    \bP_{k|k} &= (\bI-\bK_k\bH)\bP_{k|k-1}(\bI-\bK_k\bH)^\top + \bK_k\bR\bK_k^\top
\end{align}

\subsection{Key Code: Matrix Construction}
The F and Q matrices for one joint are constructed as follows:
\begin{lstlisting}[caption={Building the state transition matrix F}]
static void build_F(double* F, double dt) {
    mat_zero(F, N_STATE, N_STATE);
    double dt2 = dt*dt/2.0, dt3 = dt*dt*dt/6.0;
    // Three identical 4x4 blocks on the diagonal (x, y, z axes)
    for (int b = 0; b < 3; b++) {
        int o = b*4;
        F[(o+0)*N_STATE+(o+0)] = 1;    F[(o+0)*N_STATE+(o+1)] = dt;
        F[(o+0)*N_STATE+(o+2)] = dt2;  F[(o+0)*N_STATE+(o+3)] = dt3;
        F[(o+1)*N_STATE+(o+1)] = 1;    F[(o+1)*N_STATE+(o+2)] = dt;
        F[(o+1)*N_STATE+(o+3)] = dt*dt/2.0;
        F[(o+2)*N_STATE+(o+2)] = 1;    F[(o+2)*N_STATE+(o+3)] = dt;
        F[(o+3)*N_STATE+(o+3)] = 1;
    }
}
\end{lstlisting}

\subsection{Key Code: Predict and Update Loop}
\begin{lstlisting}[caption={Core LKF predict--update step (per joint, per frame)}]
// PREDICT
mat_mul(F, x, x_pred, N_STATE, N_STATE, 1);   // x_pred = F*x
mat_mul(F, P, FP, N_STATE, N_STATE, N_STATE);
mat_mul(FP, Ft, FPFt, N_STATE, N_STATE, N_STATE);
mat_add(FPFt, Q, P_pred, N_STATE, N_STATE);   // P_pred = F*P*Ft + Q

// INNOVATION
z[0] = noisy[k][col_offset+0];  // measured px
z[1] = noisy[k][col_offset+1];  // measured py
z[2] = noisy[k][col_offset+2];  // measured pz
mat_mul(H, x_pred, Hx, N_MEAS, N_STATE, 1);
mat_sub(z, Hx, y, N_MEAS, 1);   // y = z - H*x_pred

// INNOVATION COVARIANCE
mat_mul(H, P_pred, HP, N_MEAS, N_STATE, N_STATE);
mat_mul(HP, Ht, HPHt, N_MEAS, N_STATE, N_MEAS);
mat_add(HPHt, R, S, N_MEAS, N_MEAS);          // S = H*P*Ht + R

// KALMAN GAIN (via Cholesky -- no direct inversion)
compute_gain(P_pred, H, S, K);

// STATE UPDATE
mat_mul(K, y, Ky, N_STATE, N_MEAS, 1);
mat_add(x_pred, Ky, x, N_STATE, 1);           // x = x_pred + K*y

// COVARIANCE UPDATE (Joseph form)
mat_mul(K, H, KH, N_STATE, N_MEAS, N_STATE);
mat_sub(I12, KH, IKH, N_STATE, N_STATE);       // IKH = I - K*H
mat_mul(IKH, P_pred, tmp1, N_STATE, N_STATE, N_STATE);
mat_mul(tmp1, IKHt, tmp2, N_STATE, N_STATE, N_STATE);
mat_mul(K, R, KR, N_STATE, N_MEAS, N_MEAS);
mat_mul(KR, Kt, KRKt, N_STATE, N_MEAS, N_STATE);
mat_add(tmp2, KRKt, P, N_STATE, N_STATE);     // Joseph form
\end{lstlisting}

% ══════════════════════════════════════════════════════════════════════════════
\section{Extended Kalman Filter Implementation}

\subsection{Filter Equations}
The EKF prediction step is identical to the LKF. The update step uses the
nonlinear predicted measurement and the Jacobian:
\begin{align}
    \by_k      &= \bz_k - h(\hat{\mathbf{x}}_{k|k-1}) \\
    \mathbf{S}_k &= \bJ_k \bP_{k|k-1} \bJ_k^\top + \bR \\
    \bK_k      &= \bP_{k|k-1}\bJ_k^\top \mathbf{S}_k^{-1} \\
    \hat{\mathbf{x}}_{k|k} &= \hat{\mathbf{x}}_{k|k-1} + \bK_k \by_k \\
    \bP_{k|k}  &= (\bI - \bK_k\bJ_k)\bP_{k|k-1}(\bI - \bK_k\bJ_k)^\top + \bK_k\bR\bK_k^\top
\end{align}

\subsection{Key Code: Nonlinear Measurement and Jacobian}
\begin{lstlisting}[caption={Nonlinear measurement function h(x) and Jacobian}]
// h(x): Cartesian state -> spherical measurement
static void h_func(const double* x_state, double* z_pred) {
    double px = x_state[0], py = x_state[4], pz = x_state[8];
    double rho = sqrt(px*px + py*py) + EPS;
    double r   = sqrt(px*px + py*py + pz*pz) + EPS;
    z_pred[0]  = r;                         // range
    z_pred[1]  = manual_atan2(py, px);      // azimuth
    z_pred[2]  = manual_atan2(pz, rho);     // elevation
}

// Jacobian J (3x12): partial derivatives at predicted state
static void build_jacobian(const double* x_state, double* J) {
    mat_zero(J, N_MEAS, N_STATE);
    double px = x_state[0], py = x_state[4], pz = x_state[8];
    double rho2 = px*px + py*py + EPS;
    double rho  = sqrt(rho2);
    double r2   = px*px + py*py + pz*pz + EPS;
    double r    = sqrt(r2);
    // Row 0: dr/dstate
    J[0*N_STATE+0] = px / r;   J[0*N_STATE+4] = py / r;   J[0*N_STATE+8] = pz / r;
    // Row 1: dtheta/dstate
    J[1*N_STATE+0] = -py / rho2;    J[1*N_STATE+4] = px / rho2;
    // Row 2: dphi/dstate
    J[2*N_STATE+0] = -(px*pz)/(r2*rho);
    J[2*N_STATE+4] = -(py*pz)/(r2*rho);
    J[2*N_STATE+8] =  rho / r2;
}
\end{lstlisting}

% ══════════════════════════════════════════════════════════════════════════════
\section{Numerical Design Decisions}

\subsection{Matrix Inversion Strategy}
\label{sec:inversion}
Direct inversion of the $3\times3$ innovation covariance matrix $\mathbf{S}$ via
the formula $\bK = \bP\bH^\top\bS^{-1}$ is avoided because explicit inversion is
numerically unstable and computationally expensive when $\mathbf{S}$ is nearly
singular.

Instead, the Kalman gain is computed by \textbf{solving the linear system}
$\mathbf{S}\mathbf{K}^\top = (\bP\bH^\top)^\top$ column-by-column using
\textbf{Cholesky decomposition}. Since $\mathbf{S}$ is symmetric positive definite
by construction, Cholesky factorises $\mathbf{S} = \mathbf{L}\mathbf{L}^\top$
and then solves via forward and backward substitution:

\begin{lstlisting}[caption={Cholesky-based Kalman gain computation}]
// Instead of: K = PHt * inv(S)
// We solve: S * K[:,j] = PHt[:,j]  for each column j
for (int j = 0; j < N_MEAS; j++) {
    double ej[N_MEAS] = {0};
    ej[j] = 1.0;
    double Sinv_col[N_MEAS];
    chol_solve(S, ej, Sinv_col, N_MEAS);    // S^{-1} column j
    for (int i = 0; i < N_STATE; i++) {
        double s = 0;
        for (int k = 0; k < N_MEAS; k++) s += PHt[i*N_MEAS+k]*Sinv_col[k];
        K[i*N_MEAS+j] = s;
    }
}
\end{lstlisting}

\textbf{Why this improves stability:} Cholesky decomposition has $O(n^3/3)$ cost
compared to $O(n^3)$ for LU, and the triangular structure guarantees bounded
error amplification for well-conditioned SPD matrices.

Furthermore, the \textbf{Joseph stabilised form} is used for the covariance update:
\[
    \bP_{k|k} = (\bI-\bK\bH)\bP_{k|k-1}(\bI-\bK\bH)^\top + \bK\bR\bK^\top
\]
The standard form $\bP = (\bI-\bK\bH)\bP$ can yield a non-symmetric $\bP$ under
finite-precision arithmetic over many steps, eventually causing divergence.
The Joseph form algebraically guarantees symmetry and positive semi-definiteness.

\subsection{Joint-Wise Processing Strategy}
Rather than operating on the monolithic $276\times276$ full-body covariance matrix,
the filter is applied \textbf{independently to each of the 23 joints}. This is
valid because the joints are kinematically modelled independently (no
inter-joint coupling in $\bF$, $\bQ$, or $\bH$).

This approach reduces the dominant matrix operations from $O(276^3)$ to
$23 \times O(12^3)$, a speed improvement of approximately $276^3 / (23 \times 12^3)
\approx 1{,}900\times$ for the prediction covariance update alone.

\subsection{Memory Management}
\label{sec:memory}
The covariance matrix for one joint is $12\times12 = 144$ doubles $= 1.15$~kB.
Across 23 joints this is manageable, but several intermediate matrices (FP, FPFt,
tmp1, tmp2, etc.) are also needed.

All working matrices are \textbf{heap-allocated} using \texttt{new double[n]},
not declared as local arrays on the stack:

\begin{lstlisting}[caption={Heap allocation of large working matrices}]
// Heap-allocated: avoids stack overflow for repeated large arrays
double* F  = new double[N_STATE*N_STATE];  // 144 doubles = 1.15 kB
double* P  = new double[N_STATE*N_STATE];  // covariance
double* FP = new double[N_STATE*N_STATE];  // temporary product
// ... all intermediate matrices similarly allocated
// Freed at end of program:
delete[] F; delete[] P; delete[] FP; // etc.
\end{lstlisting}

\textbf{Justification:} Stack memory on most platforms is limited to 1--8~MB.
When running 23 joints each with $\sim$20 working matrices of $144$ doubles,
the total temporary storage needed is $\approx 23 \times 20 \times 1.15 \approx 529$~kB,
which would overflow the default stack on many systems. Heap allocation via
\texttt{new}/\texttt{delete} uses the much larger process heap (limited only by
available RAM) and is therefore both safer and more portable.

% ══════════════════════════════════════════════════════════════════════════════
\section{Manual \texttt{arctan2} Approximation}
\label{sec:atan2}

The EKF measurement model requires computing $\theta = \text{atan2}(p_y, p_x)$
and $\phi = \text{atan2}(p_z, \rho)$. The assignment explicitly forbids relying
on built-in mathematical library routines for these quantities.

\subsection{Approximation Method}
A \textbf{piecewise polynomial approximation} is used. The core identity is:
\[
    \arctan(x) \approx \frac{\pi}{4}x - x(|x|-1)\bigl(0.2732 + 0.2732|x|\bigr)
    \quad \text{for } x \in [-1, 1]
\]
This is extended to the full $\text{atan2}(y, x)$ domain via range reduction:
if $|y| > |x|$, use the identity $\arctan(y/x) = \pi/2 - \arctan(x/y)$,
then adjust the quadrant based on the signs of $x$ and $y$.

\begin{lstlisting}[caption={Manual arctan and atan2 implementation}]
// Core atan approximation on [-1, 1]
static double manual_atan(double x) {
    const double PI_4 = 0.7853981633974483;
    const double C    = 0.2732632;
    double ax = x < 0 ? -x : x;
    return PI_4 * x - x * (ax - 1.0) * (C + 0.273 * ax);
}

// Full atan2(y, x) with quadrant correction
static double manual_atan2(double y, double x) {
    const double PI = 3.14159265358979323846, PI_2 = 1.5707963267948966;
    double result;
    if (fabs(x) >= fabs(y)) {
        result = manual_atan(y / x);
        if (x < 0) result += (y >= 0) ? PI : -PI;
    } else {
        result = PI_2 - manual_atan(x / y);
        if (y < 0) result -= PI;
    }
    return result;
}
\end{lstlisting}

\subsection{Choice Justification}
\begin{itemize}[noitemsep]
    \item \textbf{Bounded error:} maximum absolute error $\approx 0.0038$ radians ($0.22^\circ$),
          which is acceptable for gait joint angle estimation.
    \item \textbf{Computational efficiency:} $O(1)$, no iterative convergence,
          suitable for the inner loop over 3{,}040 frames $\times$ 23 joints.
    \item \textbf{Correct quadrant handling:} the range-reduction and sign-based
          quadrant correction reproduce the full $(-\pi, \pi]$ output range of the
          standard \texttt{atan2}.
    \item \textbf{Singularity handling:} a small epsilon $\varepsilon = 10^{-9}$
          is added to denominators to prevent division-by-zero when the joint
          position coincides with the origin.
\end{itemize}

% ══════════════════════════════════════════════════════════════════════════════
\section{Plots and Analysis}

\subsection{State Time-Series — Pelvis Joint}
Figures~\ref{fig:pos_comparison}, \ref{fig:lkf_full}, and \ref{fig:ekf_full}
show the estimated position, velocity, acceleration, and jerk for the pelvis joint
(joint 0) across the full recording duration ($\approx 30$ seconds).

\begin{figure}[H]
    \centering
    \includegraphics[width=0.95\textwidth]{../plots/01_position_comparison.png}
    \caption{Position comparison: True vs.\ Noisy vs.\ LKF vs.\ EKF for the pelvis
             joint in $x$, $y$, $z$ axes.}
    \label{fig:pos_comparison}
\end{figure}

\begin{figure}[H]
    \centering
    \includegraphics[width=0.95\textwidth]{../plots/02_lkf_full_state.png}
    \caption{LKF full state time-series for the pelvis joint: position, velocity,
             acceleration, and jerk in all three axes.}
    \label{fig:lkf_full}
\end{figure}

\begin{figure}[H]
    \centering
    \includegraphics[width=0.95\textwidth]{../plots/03_ekf_full_state.png}
    \caption{EKF full state time-series for the pelvis joint.}
    \label{fig:ekf_full}
\end{figure}

\textbf{Observations:}
\begin{itemize}[noitemsep]
    \item Both filters substantially reduce the noise visible in the raw measurements.
    \item The position estimates closely track the true trajectory. Small phase
          delays are visible during rapid direction changes, characteristic of
          causal Kalman filtering.
    \item Velocity estimates are smooth and continuous — the filter correctly
          infers velocity from position measurements alone.
    \item Acceleration and jerk estimates show more variance than position and
          velocity, as these higher-order states are only indirectly observable.
    \item The jerk signal is bounded by the process noise model, acting as a
          built-in regulariser against non-physical oscillations.
\end{itemize}

\subsection{LKF vs.\ EKF Comparison}

\begin{figure}[H]
    \centering
    \begin{minipage}{0.55\textwidth}
        \includegraphics[width=\textwidth]{../plots/04_lkf_vs_ekf_position.png}
    \end{minipage}
    \hfill
    \begin{minipage}{0.38\textwidth}
        \includegraphics[width=\textwidth]{../plots/05_rmse_comparison.png}
    \end{minipage}
    \caption{Left: LKF vs.\ EKF position overlay. Right: RMSE bar chart by axis.}
    \label{fig:comparison}
\end{figure}

\subsection{RMSE Results}

\begin{table}[H]
    \centering
    \caption{Position RMSE (metres) for the pelvis joint}
    \begin{tabular}{@{}lccc@{}}
        \toprule
        Filter & $x$-axis & $y$-axis & $z$-axis \\
        \midrule
        LKF & -- & -- & -- \\
        EKF & -- & -- & -- \\
        \bottomrule
    \end{tabular}
    \label{tab:rmse}
    \smallskip
    \footnotesize (Values populated from \texttt{plots/rmse\_results.txt} after running the filter.)
\end{table}

\textbf{Discussion:}
\begin{itemize}[noitemsep]
    \item The LKF achieves lower position RMSE because its Cartesian measurement
          model is \textit{linear} — there is no linearisation error.
    \item The EKF introduces a small additional error from the first-order Taylor
          approximation of the spherical-to-Cartesian mapping, visible as slightly
          higher RMSE.
    \item However, the EKF is architecturally more general: it can accommodate
          radar-type sensors that natively provide range and bearing rather than
          Cartesian coordinates.
    \item Both filters significantly outperform raw noisy measurements, confirming
          that the constant-jerk model is a good fit for smooth human gait.
\end{itemize}

% ══════════════════════════════════════════════════════════════════════════════
\section{3D Full-Body Walking Simulation}

Three side-by-side 3D skeleton animations were produced using
\texttt{matplotlib.animation.FuncAnimation} with \texttt{mpl\_toolkits.mplot3d}.
Each animation shows the 23 joint positions connected by 22 body segments
corresponding to the standard mocap skeleton topology.

The three animations are:
\begin{enumerate}[noitemsep]
    \item \textbf{Measured (Noisy)} — raw sensor measurements, showing characteristic
          jitter around the true trajectory.
    \item \textbf{LKF Estimate} — smoothed trajectory from the linear filter,
          visually clean with no perceptible noise.
    \item \textbf{EKF Estimate} — smoothed trajectory from the extended filter,
          comparable smoothness to LKF with minor differences in limb extremities.
\end{enumerate}

The simulation is exported as \texttt{animation.mp4} at 20~FPS and is available
at the submission Drive link below.

\bigskip
\noindent\textbf{Drive Link:} \href{https://drive.google.com/your-link-here}
{\texttt{https://drive.google.com/your-link-here}}
\smallskip

\noindent\textit{(Replace with actual Drive link after uploading animations and CSV outputs.)}

% ══════════════════════════════════════════════════════════════════════════════
\section{Conclusion}

This milestone successfully implements both the Linear Kalman Filter and Extended
Kalman Filter in C/C++ for full-body 3D gait estimation across 23 joints and
3{,}040 frames. Key findings and contributions:

\begin{itemize}[noitemsep]
    \item A 12-state constant-jerk model accurately captures the smooth dynamics
          of human walking, enabling inference of velocity, acceleration, and jerk
          from position-only measurements.
    \item Joint-wise processing reduces computational complexity by approximately
          three orders of magnitude compared to full 276-dimensional matrix operations.
    \item Cholesky decomposition and the Joseph covariance form ensure numerical
          stability across thousands of recursive filter steps.
    \item The manual piecewise arctan2 approximation achieves sub-degree accuracy
          at $O(1)$ cost, making it suitable for the real-time embedded context for
          which these filters are designed.
    \item The LKF achieves marginally better position RMSE due to the absence of
          linearisation error, while the EKF demonstrates the generalisation of
          the framework to nonlinear measurement models.
\end{itemize}

% ──────────────────────────────────────────────────────────────────────────────
\newpage
\appendix
\section{File Structure}

\begin{verbatim}
kalman_project/
├── src/
│   ├── lkf.cpp           # Linear Kalman Filter implementation
│   └── ekf.cpp           # Extended Kalman Filter implementation
├── data/
│   ├── true.csv           # Ground truth positions
│   └── noisy.csv          # Noisy measurements
├── output/
│   ├── lkf_output.csv     # LKF filtered state (276 cols per row)
│   └── ekf_output.csv     # EKF filtered state (276 cols per row)
├── plots/
│   ├── generate_plots.py  # Python plotting script
│   ├── 01_position_comparison.png
│   ├── 02_lkf_full_state.png
│   ├── 03_ekf_full_state.png
│   ├── 04_lkf_vs_ekf_position.png
│   ├── 05_rmse_comparison.png
│   └── rmse_results.txt
├── simulation/
│   ├── animate.py         # 3D animation script
│   └── animation.mp4      # Exported animation
├── report/
│   └── report.tex         # This document
└── Makefile
\end{verbatim}

\section{Build and Run Instructions}

\begin{verbatim}
# Build both filters
make all

# Run LKF (produces output/lkf_output.csv)
./lkf

# Run EKF (produces output/ekf_output.csv)
./ekf

# Generate plots (requires Python + matplotlib + pandas + numpy)
cd plots && python generate_plots.py

# Generate 3D animation
cd simulation && python animate.py
\end{verbatim}

\end{document}
